%% kapa: exam=1 round=29 year=2018 period=1 form=A subject=economics subject_ko=경제학원론 seq=34 number=74
%% source: Q-Net 감정평가사 기출문제 (cst00309) artlSeq=5202808 file=extracted/1008/5202808/2018년 제29회 감정평가사 1차시험 기출문제/2018년 제29회 감정평가사 1차시험 기출문제/1교시 A형.pdf
%% conversion: kordoc 4.24.51 -> tools/gen1.py (auto); review=proofread-sonnet; flags=none
\begin{question}{74}
甲국은 자본이동이 완전히 자유로운 소규모 개방경제이다. 변동환율제도 하에서 甲국의 거시경제모형이 다음과 같을 때, 정책효과에 관한 설명으로 옳지 않은 것은? (단, $Y$, $M$, $r$, $e$, $p$, $r^{*}$, $p^{*}$는 각각 국민소득, 통화량, 이자율, 명목환율, 물가, 외국이자율, 외국물가이다.)
\begin{framedtext}
소비함수: $C=1000+0.5(Y-T)$ \par
투자함수: $I=1200-10000r$ \par
순수출: $NX=1000-1000\epsilon$ \par
조세: $T=1000$ \par
정부지출: $G=2000$ \par
실질환율: $\epsilon =e\frac{p}{p^{*}}$ \par
실질화폐수요: $L^{D}=40-1000r+0.01Y$ \par
실질화폐공급: $L^{S}=\frac{M}{p}$ \par
$M=5000,\ p=100,\ p^{*}=100,\ r^{*}=0.02$ \par
\end{framedtext}
\begin{choices}
  \item 정부지출을 증가시켜도 균형소득은 변하지 않는다.
  \item 조세를 감면해도 균형소득은 변하지 않는다.
  \item 통화공급을 증가시키면 균형소득은 증가한다.
  \item 확장적 재정정책을 실시하면 $e$가 상승한다.
  \item 확장적 통화정책을 실시하면 $r$이 하락한다.
\end{choices}
\end{question}
\answer{5}
